Un meteorit, de 400~kg de massa, cau sobre la Lluna amb una trajectòria perpendicular a la superfície d'aquest satèl·lit. Quan es troba a 10\,000~km de la superfície lunar, la velocitat del meteorit és de 15\,000~km/h.

\begin{apartat}{1}
Determineu el valor de la velocitat amb què el meteorit arriba a la superfície de la Lluna.
\end{apartat}

\begin{apartat}{1}
Calculeu l'energia mecànica que té el meteorit a 10\,000~km de la Lluna i la que té un cos de la mateixa massa situat en una òrbita a aquesta mateixa altura sobre la superfície de la Lluna. Indiqueu quina de les dues energies mecàniques és més gran.
\end{apartat}

\dades{%
  $G = 6,67\times 10^{-11}\ \mathrm{N\,m^2\,kg^{-2}}$\\
  $M_\mathrm{Lluna} = 7,35\times 10^{22}\ \mathrm{kg}$\\
  $R_\mathrm{Lluna} = 1,74\times 10^{6}\ \mathrm{m}$}
